\documentclass[11pt]{article}
\usepackage{amsmath,amsthm,amssymb}
\usepackage[margin=1in]{geometry}
\usepackage{hyperref}

\newtheorem{theorem}{Theorem}
\newtheorem{lemma}{Lemma}
\newtheorem{proposition}{Proposition}
\newtheorem{corollary}{Corollary}
\newtheorem*{remark}{Remark}

\title{Structural proof of $\mathrm{SL}_\gamma$: closing the column-1 SC discriminator
for $\Omega$ on $\lambda=(3,3,1,1,1)$}
\author{Clio}
\date{2026-05-20}

\newcommand{\Vp}{V_{+}}
\newcommand{\Vm}{V_{-}}
\newcommand{\Va}{V_{+}^{\alpha}}
\newcommand{\Vb}{V_{+}^{\beta}}
\newcommand{\Vg}{V_{+}^{\gamma}}

\begin{document}
\maketitle

\begin{abstract}
We close the last structural gap in the column-1 SC discriminator for the
transfer operator $\Omega$ on $\lambda=(3,3,1,1,1)$. The residual lemma
$\mathrm{SL}_\gamma$ --- that $\pi_+(R'_9 v_T) \in \Va \oplus \Vb$ for every
$\gamma$-SYT $T$ --- is proved by a single positional observation: the lower
chain $S_6 S_5 S_4 S_3$ permutes only the values $\{3,4,5,6,7\}$ and therefore
\emph{cannot move the entries $8$ and $9$}; once $8$ sits at $(5,1)$ and $9$ at
$(2,3)$, the column-1 triple is forced to come from $\{3,4,5,6,7\}$ and
\emph{automatically} realises one of the patterns $\alpha,\beta,\gamma$. There
is no room for a survivor. The factor $S_7$ then annihilates the
$\gamma$-component and $S_8$ leaves $\Va\oplus\Vb$ invariant. Combined with the
previously established $L_\alpha, L_\beta$, this completes the forward direction
of the discriminator on $\lambda=(3,3,1,1,1)$.
\end{abstract}

\section{Setup and recollection}

We work in the Hecke algebra $H_9(q)$ with Hoefsmit's semi-normal basis
$\{v_T : T\in\mathrm{SYT}(\lambda)\}$, $\lambda=(3,3,1,1,1)$, $N=120$. Write
$S_i := T_i+1$ ($i=1,\dots,8$) and
\[
R'_k := S_{k-1}S_{k-2}\cdots S_2 S_1,\qquad
\Omega := R'_3 R'_4 R'_5 R'_6 R'_7 R'_8 R'_9 .
\]
For $T\in\mathrm{SYT}(\lambda)$ let $T(r,c)$ denote the entry in row $r$, column
$c$ (both $1$-indexed). On $v_T$, the generator $T_i$ has eigenvalue $q$ if
$i,i+1$ lie in the same row of $T$ (so $S_i v_T=(q+1)v_T$), eigenvalue $-1$ if
they lie in the same column (so $S_i v_T = 0$), and otherwise
$S_i v_T = a_i v_T + b_i v_{s_i T}$, where $s_i T$ is the SYT obtained from $T$
by swapping the values $i,i+1$ (this swap fixes the cells of all other values).

Recall the $T_1$-eigenspace decomposition $V = \Vp\oplus\Vm$ with
\[
\Vp := \mathrm{span}\{v_T : T \text{ SR-at-1}\},\qquad
\Vm := \mathrm{span}\{v_T : T \text{ SC-at-1}\}=\ker S_1,
\]
where SR-at-1 means $T(1,1)=1,\,T(1,2)=2$ and SC-at-1 means $T(1,1)=1,\,T(2,1)=2$.

\paragraph{The three column-1 SC patterns (labels).}
For an SR-at-1 SYT $T$ we say
\[
\alpha:\ T(2,1)=3 \wedge T(3,1)=4;\quad
\beta:\ T(3,1)=5\wedge T(4,1)=6;\quad
\gamma:\ T(4,1)=7\wedge T(5,1)=8 .
\]
These conditions are not mutually exclusive (a SYT may be both $\alpha$ and
$\gamma$), but as \emph{sets} the spans
\[
\Va := \mathrm{span}\{v_T : T\ \alpha\},\quad
\Vb := \mathrm{span}\{v_T : T\ \beta\},\quad
\Vg := \mathrm{span}\{v_T : T\ \gamma\}
\]
are well defined subspaces of $\Vp$.

\paragraph{What is already proved.}
The companion note
\texttt{2026-05-19-col1-sc-discriminator-structural.tex} establishes
$L_\alpha$ and $L_\beta$:
\begin{itemize}
\item[$L_\alpha$:] if $T$ is $\alpha$, then $R'_k v_T \in \Vm$ for every $k\ge 4$
  (in particular $R'_k(\Va)\subseteq\Vm$);
\item[$L_\beta$:] if $T$ is $\beta$, then $\Omega v_T = 0$; moreover
  $\pi_+(R'_8 v_{T})\in\Va$ for $\beta$-SYTs (the $R'_8$-form
  $\mathrm{SL}_\beta\text{-}R8$).
\end{itemize}
The whole forward direction reduces to the single open statement below.

\begin{theorem}[$\mathrm{SL}_\gamma$]\label{thm:slg}
For every $\gamma$-SYT $T$ of $\lambda=(3,3,1,1,1)$,
\[
\pi_+(R'_9 v_T)\ \in\ \Va\oplus\Vb ,
\]
where $\pi_+\colon V\to\Vp$ is the projection along $\Vm$.
\end{theorem}

Once Theorem~\ref{thm:slg} is in hand, $L_\gamma$ (hence the discriminator)
follows verbatim from the argument already recorded in
\texttt{2026-05-19-\dots.tex} (reproduced as
Corollary~\ref{cor:Lgamma} below).

\section{The position of $8$ and $9$ is frozen}

\begin{lemma}[Frozen corners]\label{lem:frozen}
Let $T$ be a $\gamma$-SYT. Then $T(4,1)=7$, $T(5,1)=8$ and $T(2,3)=9$.
Moreover the operators $S_3,S_4,S_5,S_6$ fix the cells occupied by $8$ and $9$:
for any SYT $T'$ appearing with nonzero coefficient in
$S_6 S_5 S_4 S_3\, v_T$, one still has $T'(5,1)=8$ and $T'(2,3)=9$.
\end{lemma}

\begin{proof}
By definition of $\gamma$, $7$ is at $(4,1)$ and $8$ at $(5,1)$. The largest
entry $9$ occupies an outer corner of $\lambda$; the corners are $(1,3),(2,3),
(5,1)$. Cell $(5,1)$ is taken by $8$. If $9$ were at $(1,3)$ then column $3$
would read $T(1,3)=9 > T(2,3)$, violating column-strict increase. Hence
$9$ is at $(2,3)$.

Each $S_j=T_j+1$ acts on the basis through the value-transposition $s_j$, which
interchanges the values $j$ and $j+1$ and fixes every other cell. For
$j\in\{3,4,5,6\}$ the moved values lie in $\{3,4,5,6,7\}$, so neither $8$ nor
$9$ is ever displaced. Consequently every SYT in the support of
$S_6S_5S_4S_3v_T$ has $8$ at $(5,1)$ and $9$ at $(2,3)$.
\end{proof}

\section{No survivor can hide behind frozen corners}

This is the combinatorial heart of the proof.

\begin{lemma}[No-survivor]\label{lem:nosurv}
Let $T'$ be an SR-at-1 SYT of $\lambda=(3,3,1,1,1)$ with $T'(5,1)=8$ and
$T'(2,3)=9$. Then $T'$ is $\alpha$, $\beta$, or $\gamma$.
\end{lemma}

\begin{proof}
Fix such a $T'$. The cells $(1,1),(1,2)$ hold $1,2$ (SR-at-1); the cells
$(5,1),(2,3)$ hold $8,9$. The remaining five cells
\[
(1,3),\ (2,1),\ (2,2),\ (3,1),\ (4,1)
\]
are filled by the five values $\{3,4,5,6,7\}$. In particular the three
\emph{column-1} entries below the first row,
\[
a:=T'(2,1),\quad b:=T'(3,1),\quad c:=T'(4,1),
\]
form a strictly increasing triple $a<b<c$ chosen from $\{3,4,5,6,7\}$ (the
other two of the five values occupy $(1,3)$ and $(2,2)$).

We claim that any increasing triple $a<b<c$ from $\{3,4,5,6,7\}$ satisfies
\[
(a=3\wedge b=4)\quad\text{or}\quad(b=5\wedge c=6)\quad\text{or}\quad c=7 .
\tag{$\star$}
\]
Indeed, suppose $c\neq 7$. Being the largest of three distinct elements of
$\{3,4,5,6,7\}$, $c$ is at least $5$, so $c\in\{5,6\}$.
\begin{itemize}
\item If $c=5$ then necessarily $(a,b,c)=(3,4,5)$, giving $a=3\wedge b=4$.
\item If $c=6$ then $a<b$ are drawn from $\{3,4,5\}$, i.e.\
  $(a,b)\in\{(3,4),(3,5),(4,5)\}$. The first gives $a=3\wedge b=4$; the latter
  two give $b=5\wedge c=6$.
\end{itemize}
Thus $c\neq 7\Rightarrow (a=3\wedge b=4)\vee(b=5\wedge c=6)$, which together
with the case $c=7$ proves $(\star)$.

Translating $(\star)$ back: $a=3\wedge b=4$ means $T'(2,1)=3,T'(3,1)=4$, i.e.\
$T'$ is $\alpha$; $b=5\wedge c=6$ means $T'(3,1)=5,T'(4,1)=6$, i.e.\ $T'$ is
$\beta$; $c=7$ means $T'(4,1)=7$, and since $T'(5,1)=8$ this is exactly
$\gamma$. Hence $T'\in\alpha\cup\beta\cup\gamma$.
\end{proof}

\begin{remark}
Lemma~\ref{lem:nosurv} is the precise reason the discriminator's three strata
exhaust the residual SYTs once the corners are frozen: the column-1 ladder of a
$5$-cell first column, with top two cells $1,\ast$ and bottom cell $8$, has only
the three ``defect'' patterns $\alpha,\beta,\gamma$ available to it. A
survivor would need a column-1 triple avoiding all of
$\{3,4\text{ flush at top}\},\{5,6\text{ flush}\},\{7\text{ at }(4,1)\}$, and no
such triple exists in $\{3,4,5,6,7\}$.
\end{remark}

\begin{corollary}\label{cor:Zsupp}
For a $\gamma$-SYT $T$, $\ Z:=S_6S_5S_4S_3\,v_T\ \in\ \Va\oplus\Vb\oplus\Vg$.
\end{corollary}

\begin{proof}
The operators $S_3,S_4,S_5,S_6$ do not move the values $1,2$, so they preserve
$\Vp$ (they map SR-at-1 SYTs to SR-at-1 SYTs); thus $Z\in\Vp$. By
Lemma~\ref{lem:frozen} every SYT in the support of $Z$ has $8$ at $(5,1)$ and
$9$ at $(2,3)$, and is SR-at-1; by Lemma~\ref{lem:nosurv} each such SYT is
$\alpha,\beta$, or $\gamma$. Hence $Z\in\Va\oplus\Vb\oplus\Vg$.
\end{proof}

\section{$S_7$ collapses $\gamma$; $S_8$ stays put}

\begin{lemma}[Ladder preservation]\label{lem:ladder}
\leavevmode
\begin{enumerate}
\item $S_7\,\Vg = 0$.
\item $S_7\,\Va\subseteq\Va$ and $S_7\,\Vb\subseteq\Vb$.
\item $S_8\,\Va\subseteq\Va$ and $S_8\,\Vb\subseteq\Vb$.
\end{enumerate}
Consequently $S_7(\Va\oplus\Vb\oplus\Vg)\subseteq\Va\oplus\Vb$ and
$S_8(\Va\oplus\Vb)\subseteq\Va\oplus\Vb$.
\end{lemma}

\begin{proof}
(1) A $\gamma$-SYT has $7$ at $(4,1)$ and $8$ at $(5,1)$: the values $7,8$ are
in the same column, so $S_7 v_T = 0$.

(2)--(3) For a basis vector $v_{T'}$ we have $S_j v_{T'}\in
\mathrm{span}\{v_{T'},v_{s_j T'}\}$ (a multiple of $v_{T'}$ alone when $s_jT'$
is undefined). The transposition $s_7$ swaps the values $7,8$ and $s_8$ swaps
$8,9$; neither touches any of the values $3,4,5,6$. Therefore $s_7$ and $s_8$
leave invariant both the $\alpha$-condition (a statement about the cells of
$3,4$) and the $\beta$-condition (about the cells of $5,6$): if $T'$ is
$\alpha$ then so is $s_7T'$ (resp.\ $s_8T'$), and likewise for $\beta$. Hence
$S_7$ and $S_8$ map $\Va$ into $\Va$ and $\Vb$ into $\Vb$.

The two displayed inclusions follow: writing $w=w_\alpha+w_\beta+w_\gamma$,
$S_7w = S_7w_\alpha + S_7w_\beta + 0 \in \Va\oplus\Vb$; and $S_8$ preserves
$\Va\oplus\Vb$ by (3).
\end{proof}

\begin{remark}
In (2)--(3) we do \emph{not} claim $s_7,s_8$ preserve the $\gamma$-label
(they do not: $s_7$ turns the column-stacked $7,8$ into a column-inversion, so
$s_7T'$ is undefined on a $\gamma$-SYT, consistent with $S_7v_T=0$). We only
use that they preserve the $\alpha$- and $\beta$-labels, which is immediate
because those labels are conditions on values $\le 6$.
\end{remark}

\section{Proof of $\mathrm{SL}_\gamma$ and of $L_\gamma$}

\begin{proof}[Proof of Theorem~\ref{thm:slg}]
Let $T$ be a $\gamma$-SYT. Since $T$ is SR-at-1, $S_1 v_T=(q+1)v_T$, so
\[
R'_9 v_T = (q+1)\, S_8 S_7 S_6 S_5 S_4 S_3\, S_2\, v_T .
\]
Now $S_2 v_T = a_2 v_T + (\text{term in }\Vm)$, where the second term is a
multiple of $v_{s_2 T}$ (which is SC-at-1, as $s_2$ moves $2$ off the first row
into column~$1$) and is absent when $3$ lies in row~$1$ of $T$; the scalar
$a_2$ is $q+1$ in that case and $q/(q+1)$ otherwise --- in particular
$a_2\neq 0$. The operators $S_3,\dots,S_8$ fix the values $1,2$ and hence
preserve the splitting $V=\Vp\oplus\Vm$. Therefore
\[
\pi_+\big(R'_9 v_T\big)
= (q+1)\,\pi_+\big(S_8S_7S_6S_5S_4S_3\,S_2 v_T\big)
= (q+1)\,a_2\; S_8S_7S_6S_5S_4S_3\, v_T
= (q+1)\,a_2\; S_8 S_7 Z,
\]
with $Z=S_6S_5S_4S_3 v_T$. By Corollary~\ref{cor:Zsupp},
$Z\in\Va\oplus\Vb\oplus\Vg$. By Lemma~\ref{lem:ladder},
$S_7 Z\in\Va\oplus\Vb$ and then $S_8 S_7 Z\in\Va\oplus\Vb$. Since
$(q+1)a_2\neq 0$,
\[
\pi_+(R'_9 v_T)\in\Va\oplus\Vb . \qedhere
\]
\end{proof}

\begin{corollary}[$L_\gamma$]\label{cor:Lgamma}
If $T$ is a $\gamma$-SYT then $\Omega v_T = 0$.
\end{corollary}

\begin{proof}
Decompose $R'_9 v_T = w_+ + w_-$ with $w_-\in\Vm$. By
Theorem~\ref{thm:slg}, $w_+=\pi_+(R'_9v_T)=w_+^\alpha+w_+^\beta$ with
$w_+^\alpha\in\Va$, $w_+^\beta\in\Vb$. As each $R'_k$ has rightmost factor
$S_1$, $R'_8\Vm=0$, so $R'_8R'_9 v_T = R'_8 w_+^\alpha + R'_8 w_+^\beta$. By
$L_\alpha$ ($k=8$), $R'_8 w_+^\alpha\in\Vm$. By $\mathrm{SL}_\beta\text{-}R8$,
$\pi_+(R'_8 w_+^\beta)\in\Va$, i.e.\ $R'_8 w_+^\beta\in \Vm\oplus\Va$. Hence
\[
R'_8R'_9 v_T \in \Vm\oplus\Va .
\]
Finally apply $R'_7$: $R'_7\Vm=0$, and by $L_\alpha$ ($k=7$),
$R'_7\Va\subseteq\Vm$. Therefore $R'_7R'_8R'_9 v_T\in\Vm$, and the remaining
factors $R'_3R'_4R'_5R'_6$ (each with rightmost factor $S_1$) annihilate it:
$\Omega v_T=0$.
\end{proof}

\begin{theorem}[Discriminator, forward direction, on $(3,3,1,1,1)$]
For SR-at-1 $T\in\mathrm{SYT}(3,3,1,1,1)$: if $T$ is $\alpha$, $\beta$, or
$\gamma$ then $\Omega v_T = 0$.
\end{theorem}

\begin{proof}
$L_\alpha$, $L_\beta$, and Corollary~\ref{cor:Lgamma}.
\end{proof}

Together with the finite converse check on the $50$ SR-at-1 SYTs (recorded in
\texttt{2026-05-19-V33111-residual-diagnostic.md}), this proves the full
column-1 SC discriminator on $\lambda=(3,3,1,1,1)$ structurally.

\section{Computational verification}

All claims were verified by exact rational arithmetic at
$q\in\{2,3,5,7,11\}$ in \texttt{2026-05-20-SLgamma-verify-num.py} (parts (0)
and (A) are $q$-independent combinatorics):
\begin{itemize}
\item (0) all $8$ $\gamma$-labelled SYTs have $7@(4,1),8@(5,1),9@(2,3)$;
\item (A) the $15$ SR-at-1 SYTs with $8@(5,1),9@(2,3)$ ($6$ $\alpha$, $3$
  $\beta$, $6$ $\gamma$-only; two of the $\alpha$ are also $\gamma$) are
  \emph{all} $\alpha/\beta/\gamma$ --- no survivor, Lemma~\ref{lem:nosurv};
\item (B) $S_7 v_T=0$ for every $\gamma$-SYT, and $S_7$ preserves
  $\Va,\Vb$ --- Lemma~\ref{lem:ladder}(1)--(2);
\item (C) $S_8$ preserves $\Va,\Vb$ --- Lemma~\ref{lem:ladder}(3);
\item (D) for each of the $6$ $\gamma$-only SYTs: $Z\in\Va\oplus\Vb\oplus\Vg$,
  $Y=S_8S_7Z\in\Va\oplus\Vb$, and $\pi_+(R'_9v_T)$ is a constant multiple of
  $Y$ (identical support in $\alpha\cup\beta$).
\end{itemize}
The step-by-step support evolution
(\texttt{2026-05-20-gamma-chain-tracker.py}) shows the
$\gamma$-component dying \emph{exactly} at the $S_7$ step and the survivor- and
$\Vm$-components being absent at \emph{every} step, matching the proof.

\section{Discussion}

The proof is short because it isolates the right invariant: the
\emph{frozen corners}. Lower transpositions cannot disturb $8$ and $9$, and
with those two corners pinned the first column has no combinatorial freedom to
produce a survivor --- the column-1 ladder $\{1,\ast,\ast,\ast,8\}$ over the
alphabet $\{3,4,5,6,7\}$ realises $\alpha$, $\beta$, or $\gamma$ and nothing
else $(\star)$. This is the ``inevitability'' flavour: the discriminator's
trichotomy is forced by the shape, not engineered. The earlier attempt foundered
on trying to push $S_7$'s annihilating power down through $S_5,S_6$ via a
commutation identity; the present argument instead lets $S_3,\dots,S_6$ act
freely (they keep us inside $\Va\oplus\Vb\oplus\Vg$ for free) and only invokes
$S_7$ at the top of the chain, where its action on $\gamma$ is the trivial
column-collapse $S_7 v_T=0$.

\paragraph{Toward generalisation.} The mechanism suggests the discriminator on a
general SR-at-1 stratum: freeze the largest corner entries, observe that the
column-1 ladder over the residual alphabet admits only the SRD ``defect''
patterns at the column-1 promotion indices, then collapse them top-down with the
matching $S_i$. This is the route to the conjectured uniform statement
(Open puzzle E.1 of the diagnostic).

\section*{Files}
\begin{itemize}
\item Proof: this file.
\item Verification: \texttt{\textasciitilde/projects/scratch/2026-05-20-SLgamma-verify.py}
\item Mechanism trace: \texttt{\textasciitilde/projects/scratch/2026-05-20-gamma-chain-tracker.py}
\item Prior lemmas: \texttt{\textasciitilde/projects/proofs/2026-05-19-col1-sc-discriminator-structural.tex}
\end{itemize}

\end{document}
